\documentclass[a4paper,10pt]{report}\input{../0.Préambule/Préambule_cours_prof}\begin{document}

\pagestyle{empty}
\newpage
\noindent NOM, Prénom et classe : \dotfill

\section*{Les triangles 2 \hspace{8cm} Sujet A}
\addcontentsline{toc}{section}{Les triangles 2 }

\noindent \textit{Il sera tenu compte dans la notation de la propreté ainsi que de la justification apportée à chacune des réponses.}

\noindent \textit{Le barème est donné à titre indicatif. Il pourra être modifié ultérieurement.}

\noindent \textit{L'usage de la calculatrice est \textbf{interdite}.} \hfill \textit{Durée : 30 minutes}

\paragraph{Exercice 1}: \hfill \textbf{4 points}

Les tracés seront soignés, les points nommés et les traits de construction laissés.
\begin{enumerate}
\item Construire un triangle $LMN$ tel que $LM = 8$cm, $MN = 5$cm et $\widehat{LMN} =120$\degre .
\item Construire un triangle $PQR$ tel que $PQ = 7$cm, $\widehat{PQR} = 40$\degre et $\widehat{QPR} =50$\degre .
\end{enumerate}

\vspace{8cm}

\paragraph{Exercice 2}: \hfill \textbf{6 points}

\begin{enumerate}
\item Préciser la nature d'un triangle possédant un angle mesurant $24$\degre et un autre angle mesurant $66$\degre . Justifie.
\item Préciser la nature d'un triangle possédant un angle mesurant $24$\degre et un autre angle mesurant $60$\degre . Justifie.
\end{enumerate}

\noindent\grandscarreaux{\linewidth}{8cm}

\newpage
\paragraph{Exercice 3}: \hfill \textbf{5 points}

Recopier chacune des phrases sur ta feuille en complétant les pointillés.
\begin{enumerate}
\item Si $BUS$ est un triangle isocèle en $U$ alors les deux côtés [\dots ~\dots ~] et [\dots ~\dots ~] issus du sommet $U$ sont de même \dots ~\dots ~ et les deux angles à sa base $\widehat{\dots ~\dots ~}$ et $\displaystyle \widehat{\dots ~\dots ~}$ ont même \dots ~\dots ~ .
\item Si un triangle $CAR$ est rectangle en $C$ alors l'angle $\widehat{ACR}$ est un angle \dots ~\dots ~ et la somme des angles $\widehat{\dots ~\dots ~}$ et $\widehat{\dots ~\dots ~}$ vaut \dots ~\dots ~ .
\end{enumerate}

\noindent\grandscarreaux{\linewidth}{8cm}

\medskip
\paragraph{Exercice 4}: \hfill \textbf{5 points}

\noindent Observe les codages de la figure suivante puis calcule la mesure de l'angle $\widehat{IOS}$. On Justifiera la réponse en faisant apparaitre les calculs effectués

\noindent \begin{minipage}{8cm}
\noindent\grandscarreaux{\linewidth}{5.6cm}
\end{minipage}
\hspace{5mm}
\begin{minipage}{8cm}
\begin{center}
\begin{tikzpicture}[line cap=round,line join=round,>=triangle 45,x=1.0cm,y=1.0cm,scale=0.8]
\clip(0.2,0.2) rectangle (10.2,6.2);
\draw [shift={(9.42,2.48)},line width=0.8pt,color=red,fill=red,fill opacity=0.1] (0,0) -- (159.28023015448653:0.6) arc (159.28023015448653:191.28023015448653:0.6) -- cycle;
\draw[line width=0.8pt,color=red,fill=red,fill opacity=0.1] (2.4160682250247367,1.0829893494658505) -- (2.3330788755588863,1.4990575744905872) -- (1.9170106505341495,1.4160682250247367) -- (2.,1.) -- cycle; 
\draw [line width=1.2pt] (2.,1.)-- (9.42,2.48);
\draw [line width=1.2pt] (2.,1.)-- (5.761877543128084,3.8637319351247728);
\draw [line width=1.2pt] (3.7885400817040114,2.4872056649969556) -- (3.9096828590940165,2.328069159171814);
\draw [line width=1.2pt] (3.8521946840340684,2.5356627759529577) -- (3.9733374614240735,2.3765262701278163);
\draw [line width=1.2pt] (9.42,2.48)-- (1.0751933591741956,5.63653059116721);
\draw [line width=1.2pt] (2.,1.)-- (1.0751933591741956,5.63653059116721);
\draw [line width=1.2pt] (1.4473527693731592,3.2594772298789185) -- (1.6434892116073194,3.2985987843946276);
\draw [line width=1.2pt] (1.431704147566876,3.3379318067725827) -- (1.6278405898010362,3.377053361288292);
\begin{scriptsize}
\draw [color=black] (2.,1.)-- ++(-3.0pt,0 pt) -- ++(6.0pt,0 pt) ++(-3.0pt,-3.0pt) -- ++(0 pt,6.0pt);
\draw[color=black] (1.62,0.67) node {\large{$O$}};
\draw [color=black] (9.42,2.48)-- ++(-3.0pt,0 pt) -- ++(6.0pt,0 pt) ++(-3.0pt,-3.0pt) -- ++(0 pt,6.0pt);
\draw[color=black] (9.84,2.87) node {\large{$L$}};
\draw[color=red] (8,2.67) node {\large{$36\textrm{\degre}$}};
\draw [color=black] (1.0751933591741956,5.63653059116721)-- ++(-3.0pt,0 pt) -- ++(6.0pt,0 pt) ++(-3.0pt,-3.0pt) -- ++(0 pt,6.0pt);
\draw[color=black] (0.44,5.93) node {\large{$I$}};
\draw [color=black] (5.761877543128084,3.8637319351247728)-- ++(-3.0pt,0 pt) -- ++(6.0pt,0 pt) ++(-3.0pt,-3.0pt) -- ++(0 pt,6.0pt);
\draw[color=black] (5.9,4.27) node {\large{$S$}};
\end{scriptsize}
\end{tikzpicture}
\end{center}
\end{minipage}

\vspace{-0.6mm}
\noindent\grandscarreaux{\linewidth}{5.6cm}

\newpage
\noindent NOM, Prénom et classe : \dotfill

\section*{Les triangles 2 \hspace{8cm} Sujet B}
\addcontentsline{toc}{section}{Les triangles 2}

\noindent \textit{Il sera tenu compte dans la notation de la propreté ainsi que de la justification apportée à chacune des réponses.}

\noindent \textit{Le barème est donné à titre indicatif. Il pourra être modifié ultérieurement.}

\noindent \textit{L'usage de la calculatrice est \textbf{interdite}.} \hfill \textit{Durée : 30 minutes}

\paragraph{Exercice 1}: \hfill \textbf{4 points}

Les tracés seront soignés, les points nommés et les traits de construction laissés.
\begin{enumerate}
\item Construire un triangle $LMN$ tel que $LM = 5$cm, $MN = 8$cm et $\widehat{LMN} =110$\degre .
\item Construire un triangle $PQR$ tel que $PQ = 7$cm, $\widehat{PQR} = 50$\degre et $\widehat{QPR} =40$\degre .
\end{enumerate}

\vspace{8cm}

\paragraph{Exercice 2}: \hfill \textbf{6 points}

\begin{enumerate}
\item Préciser la nature d'un triangle possédant un angle mesurant $24$\degre et un autre angle mesurant $68$\degre . Justifie.
\item Préciser la nature d'un triangle possédant un angle mesurant $48$\degre et un autre angle mesurant $84$\degre . Justifie.
\end{enumerate}

\noindent\grandscarreaux{\linewidth}{8cm}

\newpage
\paragraph{Exercice 3}: \hfill \textbf{5 points}

Recopier chacune des phrases sur ta feuille en complétant les pointillés.
\begin{enumerate}
\item Si un triangle $BUS$ est rectangle en $B$ alors l'angle $\widehat{UBS}$ est un angle \dots ~\dots ~ et la somme des angles $\widehat{\dots ~\dots ~}$ et $\widehat{\dots ~\dots ~}$ vaut \dots ~\dots ~ .
\item Si $CAR$ est un triangle isocèle en $A$ alors les deux côtés [\dots ~\dots ~] et [\dots ~\dots ~] issus du sommet $U$ sont de même \dots ~\dots ~ et les deux angles à sa base $\widehat{\dots ~\dots ~}$ et $\displaystyle \widehat{\dots ~\dots ~}$ ont même \dots ~\dots ~ .

\end{enumerate}

\noindent\grandscarreaux{\linewidth}{8cm}

\medskip
\paragraph{Exercice 4}: \hfill \textbf{5 points}

\noindent Observe les codages de la figure suivante puis calcule la mesure de l'angle $\widehat{IOS}$. On Justifiera la réponse en faisant apparaitre les calculs effectués

\noindent \begin{minipage}{8cm}
\noindent\grandscarreaux{\linewidth}{5.6cm}
\end{minipage}
\hspace{5mm}
\begin{minipage}{8cm}
\begin{center}
\begin{tikzpicture}[line cap=round,line join=round,>=triangle 45,x=1.0cm,y=1.0cm,scale=0.8]
\clip(0.2,0.2) rectangle (10.2,6.2);
\draw [shift={(9.42,2.48)},line width=0.8pt,color=red,fill=red,fill opacity=0.1] (0,0) -- (159.28023015448653:0.6) arc (159.28023015448653:191.28023015448653:0.6) -- cycle;
\draw[line width=0.8pt,color=red,fill=red,fill opacity=0.1] (2.4160682250247367,1.0829893494658505) -- (2.3330788755588863,1.4990575744905872) -- (1.9170106505341495,1.4160682250247367) -- (2.,1.) -- cycle; 
\draw [line width=1.2pt] (2.,1.)-- (9.42,2.48);
\draw [line width=1.2pt] (2.,1.)-- (5.761877543128084,3.8637319351247728);
\draw [line width=1.2pt] (3.7885400817040114,2.4872056649969556) -- (3.9096828590940165,2.328069159171814);
\draw [line width=1.2pt] (3.8521946840340684,2.5356627759529577) -- (3.9733374614240735,2.3765262701278163);
\draw [line width=1.2pt] (9.42,2.48)-- (1.0751933591741956,5.63653059116721);
\draw [line width=1.2pt] (2.,1.)-- (1.0751933591741956,5.63653059116721);
\draw [line width=1.2pt] (1.4473527693731592,3.2594772298789185) -- (1.6434892116073194,3.2985987843946276);
\draw [line width=1.2pt] (1.431704147566876,3.3379318067725827) -- (1.6278405898010362,3.377053361288292);
\begin{scriptsize}
\draw [color=black] (2.,1.)-- ++(-3.0pt,0 pt) -- ++(6.0pt,0 pt) ++(-3.0pt,-3.0pt) -- ++(0 pt,6.0pt);
\draw[color=black] (1.62,0.67) node {\large{$O$}};
\draw [color=black] (9.42,2.48)-- ++(-3.0pt,0 pt) -- ++(6.0pt,0 pt) ++(-3.0pt,-3.0pt) -- ++(0 pt,6.0pt);
\draw[color=black] (9.84,2.87) node {\large{$L$}};
\draw[color=red] (8,2.67) node {\large{$42\textrm{\degre}$}};
\draw [color=black] (1.0751933591741956,5.63653059116721)-- ++(-3.0pt,0 pt) -- ++(6.0pt,0 pt) ++(-3.0pt,-3.0pt) -- ++(0 pt,6.0pt);
\draw[color=black] (0.44,5.93) node {\large{$I$}};
\draw [color=black] (5.761877543128084,3.8637319351247728)-- ++(-3.0pt,0 pt) -- ++(6.0pt,0 pt) ++(-3.0pt,-3.0pt) -- ++(0 pt,6.0pt);
\draw[color=black] (5.9,4.27) node {\large{$S$}};
\end{scriptsize}
\end{tikzpicture}
\end{center}
\end{minipage}

\vspace{-0.6mm}
\noindent\grandscarreaux{\linewidth}{5.6cm}
\end{document}